% TOP_PROBLEM: {"problemId": "problem.maximum-number-of-k-sets-and-complexity-of-k-levels", "canonicalTitle": "Maximum Number of k-sets and Complexity of k-levels", "releaseRank": 479, "primaryDomain": "combinatorics_discrete_geometry", "primaryDomainLabel": "Combinatorics & discrete geometry", "displayStatus": "open", "releaseStatus": "open"}
% !TeX program = lualatex
% ATTEMPT_STATUS: unresolved
\documentclass[11pt]{article}
\usepackage[margin=25mm]{geometry}
\usepackage{fontspec}
\setmainfont{Latin Modern Roman}
\usepackage{amsmath,amssymb,mathtools}
\usepackage{unicode-math}
\setmathfont{Latin Modern Math}
\usepackage{xurl,hyperref}
\hypersetup{colorlinks=true,urlcolor=blue}
\setcounter{secnumdepth}{0}
\setlength{\parindent}{0pt}
\setlength{\parskip}{5pt}
\setlength{\emergencystretch}{3em}
\begin{document}
\section*{\texorpdfstring{Maximum Number of k-sets and Complexity of k-levels}{Maximum Number of k-sets and Complexity of k-levels}}\label{479}
\subsection{Definitions and mathematical statement}
For an $n$-point set $P\subset{\mathbb{R}}^2$, a $k$-set is a subset $Q\subset P$ with $|Q|=k$ that can be strictly separated from $P\setminus Q$ by a line. Let
\[
 K(n,k)=\max_{|P|=n}\#\{Q\subset P:Q\text{ is a }k\text{-set}
                                      \},\qquad1\le k<n.
\]
Determine its extremal growth as a function of $n$ and $k$, with the general-position convention of the cited problem. In a simple arrangement of nonvertical lines, the $k$-level consists of points on the lines with exactly $k$ lines strictly below, with closures at vertices. Point--line duality relates the two extremal problems, with a possible index shift depending on convention.

The planar case corresponds to line arrangements. Arrangements of hyperplanes in higher dimensions introduce an additional dimension parameter; they are related generalizations, not literally the same planar function.
\par\smallskip\textit{Formulation and concept references:} \href{https://topp.openproblem.net/p7}{[S1]}.

\subsection{Short English statement}
What is the largest possible number of k-point subsets cut off by a line from an n-point planar set, across all values of n and k?
\subsection{Research attempt}
\paragraph{Scope, literature check, and conventions.}
This attempt by GPT-6 Astra Ultra was prepared on 27 September 2026.
Write $K_P(k)$ for the number of $k$-sets of one particular planar
configuration $P$. The problem asks for the maximum over configurations,
not the average over random points or over $k$. Complementation gives
$K_P(k)=K_P(n-k)$, so the interesting parameter can be restricted to
$1\leq k\leq n/2$.

The catalogue source [S1] retains this as an open problem. Dey's primary
paper [E1] gives the upper bound $O(nk^{1/3})$ in that range, and the
publisher abstract of T\'oth's construction [E2] gives a lower bound
$n\exp(\Omega(\sqrt{\log k}))$. Neither bound is reproved here or
advertised as a new result. A consulted 2024 paper [E3] still identifies
the planar $O(n^{4/3})$ halving bound as the available general upper
scale. Bounds for a random model or for a different dimension would not
close the gap in this planar worst-case question.

Assume no three points are collinear. For sweep bookkeeping it is convenient
also to assume distinct first coordinates and distinct directions of
segments joining pairs of points. These additional generic conditions
can be obtained by an arbitrarily small perturbation. Every strictly
separated subset in the original finite configuration has a separator with
a positive margin, so a sufficiently small simultaneous perturbation
preserves every existing $k$-set. Consequently using the stronger generic
conditions does not reduce the extremal supremum. The exact identities
below are first proved under those conditions, avoiding simultaneous events.

The strategy is to encode the whole sequence $K_P(k)$ by directed edges,
derive constraints from random convex hulls, and determine whether those
constraints can control an individual level. This yields exact reductions
and a complete elementary shallow-level estimate. An explicit nongeometric
profile then shows why those upper constraints alone cannot recover even
the known subquadratic halving bound.

\paragraph{Lemma 479.1: count changes of the lowest $k$ projections.}
For distinct $a,b\in P$, let $r(a,b)$ be the number of points strictly
to the right of the directed line from $a$ to $b$, and put
\[
 E_j(P)=\#\{(a,b)\in P^2:a\ne b,\ r(a,b)=j\},
                  \qquad 0\leq j\leq n-2.
\]
Then
\begin{equation}\label{a479:edgeidentity}
                               K_P(k)=E_{k-1}(P).
\end{equation}

\emph{Proof.} A subset $Q$ is strictly separable exactly when there is
a nonzero vector $u$ such that
\[
                      u\cdot q<u\cdot p
                \qquad(q\in Q,\ p\in P\setminus Q).
\]
For such a $u$, $Q$ is the set of the $k$ smallest projections onto $u$.
The set of feasible $u$ is an intersection of open homogeneous halfplanes,
so it is an open convex cone. It contains no antipodal pair, since reversing
the direction reverses every strict comparison. Its intersection with the
unit circle is therefore one connected open arc, if nonempty: positive
linear combinations join any two feasible directions along the intervening
short circular arc. Thus a given $Q$ appears only once as the direction
rotates around the circle.

Away from directions perpendicular to a segment joining two points, the
projection order is constant. At such an event exactly two points tie,
and the lowest-$k$ subset changes if and only if they occupy positions
$k$ and $k+1$. For the normal obtained by rotating $b-a$ counterclockwise
through $\pi/2$, the points with smaller projection than $a,b$ are
exactly those to the right of the directed line $ab$. Hence the event
changes the subset exactly when $r(a,b)=k-1$. Each ordered pair represents
one of the two antipodal events for its underlying segment.

The open arcs for the distinct $k$-sets form a cyclic succession, and
the number of changes equals the number of arcs. There is more than one
arc because a nonempty proper subset cannot be lowest in both opposite
directions. Counting the changes proves \eqref{a479:edgeidentity}.
\hfill$\square$

This proof checks why counting appearances during a sweep does not
overcount a subset appearing in several disconnected direction intervals.
Convexity of its feasible cone is the missing justification in a bare
``rotate a line'' argument.

\paragraph{Corollaries: exact total mass, symmetry, and a benchmark family.}
Reversing a directed edge interchanges its two sides, giving
\begin{equation}\label{a479:symmetry}
                    E_j(P)=E_{n-2-j}(P).
\end{equation}
Every ordered pair belongs to exactly one of the counts, so
\begin{equation}\label{a479:total}
       \sum_{k=1}^{n-1}K_P(k)=\sum_{j=0}^{n-2}E_j(P)=n(n-1).
\end{equation}
The average over $k$ is exactly $n$ for every generic configuration.
This does not say that its largest level is $O(n)$, and it does not allow
interchanging the maximum over $P$ with the sum over $k$.

If $P$ consists of the vertices of a strictly convex polygon, then
\begin{equation}\label{a479:convex}
                              K_P(k)=n\qquad(1\leq k<n).
\end{equation}
Indeed a line avoiding the vertices cuts off a consecutive block of vertices
in cyclic order: its intersection with the polygon boundary has at most
two points. Conversely any nonempty proper consecutive block is separated
by a chord joining interiors of the two boundary edges that delimit the
block. Convexity puts the intervening chord inside the polygon and the two
vertex arcs on opposite sides. There are precisely $n$ such blocks of
size $k$. Thus $K(n,k)\geq n$, and the constant profile $E_j=n$
realizes the average identity simultaneously at every level. Convexity is
an easy family, not a candidate extremal construction for large $k$.

When $n=2m$ and $k=m$, an undirected halving segment contributes both
orientations to $E_{m-1}$. Consequently the number of halving \emph{subsets}
is twice the number of halving \emph{segments}. This factor is retained
in the comparison with dual arrangements.

\paragraph{Exact dual indexing rather than a hidden shift.}
For a point $p=(a_p,b_p)$ use the dual nonvertical line
\[
                              \ell_p(t)=a_pt-b_p.
\]
Distinct $a_p$ and no three collinear primal points give a simple line
arrangement. Let $V_j$ denote the number of intersections of two dual
lines with exactly $j$ \emph{other} lines strictly below their common
height. Then
\begin{equation}\label{a479:duality}
                      K_P(k)=V_{k-1}+V_{n-k-1}.
\end{equation}
To prove this, directions of the form $(t,-1)$ cover one open semicircle
of projection directions. The $k$ smallest projections change at exactly
the crossings with $k-1$ lines below. The opposite semicircle corresponds
to the $k$ largest values of $a_pt-b_p$, equivalently the complement of
the $n-k$ smallest; its changes are counted by $V_{n-k-1}$. The two
omitted horizontal normal directions have no tie because the $a_p$ are
distinct, so there is no uncounted boundary event.

There is a further convention issue for the level as defined in the root
statement. Let $L_j$ be the closure of the open line portions with exactly
$j$ lines below. At a crossing with $q$ other lines below, the lower
two incident portions have $q$ lines below and the upper two have $q+1$.
It follows that the number of vertices on $L_j$ is
\begin{equation}\label{a479:levelvertices}
                             V_{j-1}+V_j,
             \qquad V_{-1}=V_{n-1}=0.
\end{equation}
Thus the exact identity for one fixed $k$-set count is
\eqref{a479:duality}, while a level's vertex count has
\eqref{a479:levelvertices}. They are closely related extremal quantities,
but one should not silently replace one by the other on an individual
configuration. In the halving case \eqref{a479:duality} becomes
$K_P(m)=2V_{m-1}$, illustrating the distinction explicitly.

\paragraph{Lemma 479.2: a hull-sampling generating identity.}
Independently retain every point with probability $p\in(0,1]$, obtaining
$R\subseteq P$. Let $h(R)$ be the number of directed hull edges, with
$h(R)=0$ for $|R|\leq1$, $h(R)=2$ for $|R|=2$, and the usual number
of convex-hull vertices for larger noncollinear samples. Then
\begin{equation}\label{a479:sampling}
                   \mathbb E h(R)
                   =p^2\sum_{j=0}^{n-2}E_j(P)(1-p)^j
                   \leq np.
\end{equation}

\emph{Proof.} An ordered pair $(a,b)$ becomes a hull edge with the
interior of the sampled hull on its left exactly when $a,b$ are both
retained and every original point to its right is omitted. Its probability
is $p^2(1-p)^{r(a,b)}$. Points on its left are unrestricted; no other
point is on the line. Linearity of expectation gives the equality,
including the two-point sample where both orientations count. Finally
$h(R)\leq|R|$ in every case, and $\mathbb E|R|=np$.
\hfill$\square$

For $k\geq1$, take $p=1/(k+1)$. For $0\leq j\leq k-1$,
\[
            (1-p)^j\geq(k/(k+1))^{k-1}\geq e^{-1},
\]
where the last inequality follows from
$\log(1+1/k)\leq1/k$. Since all terms are nonnegative,
\begin{equation}\label{a479:shallow}
       \sum_{j=0}^{k-1}E_j(P)\leq e n(k+1),\qquad
                            K_P(k)\leq e n(k+1).
\end{equation}
This is a complete elementary shallow-edge estimate. It is substantially
weaker than the known $O(nk^{1/3})$ bound for one level and does not
improve the current frontier.

Optimizing the sampling inequality for a single term confirms the same
limitation. For $k\geq2$, it gives
\[
                      E_{k-1}\leq\frac{n}{p(1-p)^{k-1}}.
\]
The denominator is largest at $p=1/k$, because the derivative of its
logarithm is $1/p-(k-1)/(1-p)$. The result is asymptotic to $enk$.
Changing only the sample probability does not manufacture the missing
factor $k^{2/3}$.

\paragraph{A profile that defeats the proposed upper-constraint argument.}
Could using \emph{all} probabilities $p$, the total mass identity, and
symmetry strengthen this method enough? These particular constraints
alone cannot even force a subquadratic middle level. The following
construction is a formal numerical profile, deliberately not claimed to
come from a point set.

For even $n\geq48$, put $r=(n-2)/2$ and define
\begin{equation}\label{a479:fakeprofile}
 \widetilde E_0=\widetilde E_{n-2}=3,\qquad
 \widetilde E_r=n(n-1)-6,\qquad
 \widetilde E_j=0\quad\text{at all other indices}.
\end{equation}
It is nonnegative, integral, symmetric, has total mass $n(n-1)$, and
has the minimum three supporting edges at each extreme. Nevertheless
its middle entry is quadratic. It satisfies every \emph{upper} inequality
in \eqref{a479:sampling}:
\begin{equation}\label{a479:fakeupper}
                p\sum_j\widetilde E_j(1-p)^j\leq n
                              \qquad(0\leq p\leq1).
\end{equation}

\emph{Verification.} Write $q=1-p$ and $M=n(n-1)-6$. The left side is
$3p(1+q^{2r})+Mpq^r$. Its first term is at most six. The maximum of
$pq^r$ is $(r/(r+1))^r/(r+1)$. For $r\geq2$,
\[
                  (r/(r+1))^r\leq4/9.
\]
For example, monotonicity follows by expanding
$(1+1/r)^r=\sum_{j=0}^r(j!)^{-1}\prod_{h=0}^{j-1}(1-h/r)$:
each existing factor increases with $r$ and extra terms are positive.
The value at $r=2$ is $9/4$. Since $r+1=n/2$,
\[
 3p(1+q^{2r})+Mpq^r
 \leq6+\frac{8M}{9n}
 \leq6+\frac89(n-1)\leq n
\]
for even $n\geq48$. This proves \eqref{a479:fakeupper}.

There is no contradiction with Dey's theorem: the profile is not asserted
to be geometrically realizable. Indeed random hulls also obey lower
constraints not used in this upper-bound argument, and geometric
realizability imposes much more. The example proves the narrower and
useful statement that the total identity, reversal symmetry, the extreme
count, and the entire family of the displayed upper sampling inequalities
do not imply a subquadratic middle bound. An argument depending only on
those data cannot solve the extremal problem.

\paragraph{Where the geometric information was lost.}
The random-hull calculation sees the distribution of numbers of points on
one side of one directed edge. It does not track how different edges cross,
how their endpoints interleave, or how rank transitions belonging to
different levels fit into one arrangement. The formal profile exploits
exactly that loss. One cannot insert it as a point-set construction, nor
exclude it by appealing to sampling inequalities that it already satisfies.

A natural next route is to draw the directed edges with $k-1$ points on
their right and use their crossings. This is the kind of additional
geometric input in the established approach of [E1]. An attempted shortcut
was to assume that each point can be charged only a bounded number of
rank-$k$ transitions, because it changes its membership in a lowest-$k$
set only when it ties with another point. But a point has $n-1$ possible
tie partners, and the rank of the tie varies with direction; no constant
bound follows. Counting each pair once only recovers
\eqref{a479:total}, which allows concentration on a single level.

The first missing statement in the present proof track is a bound that
couples transitions at an individual rank to their geometric crossings
or to neighboring ranks strongly enough to exclude that concentration.
Simply restating the known $O(nk^{1/3})$ theorem as such a statement
would not improve it or determine the true scale. No stronger coupling
inequality is established here.

\paragraph{Exact finite verification of both interpretations.}
The companion script checks 35 configurations: one convex configuration
and four fixed-seed generic integer configurations at each size
$3\leq n\leq9$. The convex examples use
$(2^i,4^i)$, whose pair slopes $2^i+2^j$ are distinct. Random candidates
are rejected if any three points are collinear, any first coordinates
coincide, or any two segments have parallel directions.

For every configuration the script computes all directed side counts by
integer determinants. Independently, it samples projection directions on
both sides of every critical event, enumerates the resulting lowest-$k$
subsets, and deduplicates them as exact index sets. Directions are made
with integer vectors of the form $M u\pm d$, where $u$ is perpendicular
to a segment vector $d$ and $M$ exceeds every relevant dot-product
perturbation. This keeps all nontied comparisons unchanged; no floating-point
angle or tolerance is used. The resulting counts agree with
\eqref{a479:edgeidentity}, symmetry, and the total mass identity.

The dual arrangement intersections are computed with rational arithmetic,
and \eqref{a479:duality} is checked separately. For configurations with
at most eight points, every sampled subset is enumerated and its directed
hull edges counted; the exact expectation at $p=1/2$ agrees with
\eqref{a479:sampling}. The artificial profile at $n=48$ is also tested
at 999 rational probabilities; its validity for all $p$ rests on the
analytic proof above, not on that finite grid. Source, script, and result
records are retained in
\nolinkurl{_Local/top_problems/continuation_488_477_2026_09_27/algorithms/}.

\paragraph{Outcome and explicit remaining gap.}
\textbf{UNRESOLVED.} The work proves the directed-edge count, the exact
dual indexing formulas, a convex-position family, and an elementary
$O(nk)$ shallow-edge estimate. It also constructs a rigorous obstruction
to extracting a sharp individual-level bound from the upper hull-sampling
inequalities alone. These are verified reductions and failure diagnostics,
not new matching extremal bounds.

To continue, one could add a crossing statistic to the profile, seek an
inequality relating it to the local rank sequence at each endpoint, and
test that inequality on both the convex examples and known many-$k$-set
constructions. Any candidate construction must supply actual planar
coordinates or a realizability proof; a consistent list of side counts
does not suffice. Neither the gap between the source bounds nor the full
dependence on $n$ and $k$ has been resolved by this attempt.

\subsection{Sources}
\begin{itemize}
\item[S1]Formal statement, Problem 7: Statement.
\url{https://topp.openproblem.net/p7}.
\textit{Formulation source. Consulted 24 September 2026: Planar k-set definition and the line-arrangement comparison.}
\item[E1]Tamal K. Dey, Improved Bounds for Planar $k$-Sets and Related
Problems, Discrete \& Computational Geometry 19 (1998), 373--382.
\url{https://link.springer.com/article/10.1007/PL00009354}.
\textit{Primary upper-bound paper; publisher record and author-uploaded
text consulted 27 September 2026.}
\item[E2]G\'eza T\'oth, Point Sets with Many $k$-Sets,
Discrete \& Computational Geometry 26 (2001), 187--194.
\url{https://link.springer.com/article/10.1007/s004540010022}.
\textit{Publisher abstract verifies the stated lower bound for $n\geq2k$;
the full construction is not claimed independently checked here.
Consulted 27 September 2026.}
\item[E3]Ranita Biswas, Sebastiano Cultrera di Montesano, Herbert
Edelsbrunner, and Morteza Saghafian, Depth in arrangements:
Dehn--Sommerville--Euler relations with applications,
Journal of Applied and Computational Topology 8 (2024), 557--578.
\url{https://link.springer.com/article/10.1007/s41468-024-00173-w}.
\textit{Recent primary paper's frontier comparison; consulted
27 September 2026.}
\end{itemize}
\end{document}
